\subsection{What invariant forms do not control}

\begin{proposition}[B.4.29; Proved; full-rank periods do not bound an affine comparison]
\label{push:B429}
There is a proper semistable family of two-dimensional
complex tori over a disc, with smooth total threefold
and rank-two unipotent monodromy on $H^1$, and a
fibre-preserving automorphism over the punctured disc
with the following properties. It acts identically
on all relative invariant one- and two-forms and
on the Gauss--Manin variation. Their coherent
extensions supplied by the proper family therefore
give the same lattices under this identification.
An absolute holomorphic canonical form extending
across the centre is also unchanged. Nevertheless,
the automorphism has no meromorphic extension to
the total family.

Thus these invariant-form and connection data do
not alone control an affine origin or a marked
total-space comparison, even in the presence of
a proper smooth semistable cap. This statement
does not refute a separate extension theorem
for the difference of two sections which both
extend holomorphically into the same cap.
\end{proposition}

\begin{proof}
\emph{A proper full-rank example.}
Let $\mathcal E\to\Delta_t$ be the standard Tate
family with ordinary fibre
\[
 E_t=\C^*/t^{\Z}\qquad(t\ne0)
\]
and reduced nodal rational central fibre. Its
proper analytic model and node chart $xy=t$
are the explicit example in
\citet[Sections 1.1.1--1.1.3, pp.~73--74]{KKN2008Root}.
Form $\mathcal E\times_{\Delta}\mathcal E$.
Its only total-space singularity is the product
of the two nodes, with local equations
\[
 xy=t,\qquad uv=t,\qquad xy=uv.
\]
Blow up the ideal $(x,u)$ in this conifold and
take its strict-transform charts. They are
\[
 u=xA,\quad y=Av,\quad t=xAv
\]
with smooth coordinates $(x,A,v)$, and
\[
 x=uB,\quad v=By,\quad t=uBy
\]
with smooth coordinates $(u,B,y)$. This is the
standard proper small resolution, an isomorphism
away from the conifold point. Glue it to the
identity on the complement of that point. The
resulting $p:\mathcal Y\to\Delta$ is proper,
has smooth total threefold, and has reduced
normal-crossing centre: its displayed equations
are products of distinct coordinates. The
resolution does not alter the ordinary fibres
$E_t\times E_t$.

For each elliptic factor choose period cycles
$w_i,u_i$ with $u_i\mapsto u_i+w_i$ around
$t=0$. On the direct sum the logarithm of
monodromy has rank two and square zero.
The same rank holds on dual $H^1$. Thus this
example has the full-rank weight-one
unipotent degeneration at issue, rather than
a constant elliptic factor.

\emph{The nonextending translation.}
On the punctured disc use multiplicative
coordinates $(Z,W)$ for the two factors and put
\[
 T_h(Z,W,t)=\bigl(e^{2\pi i h(t)}Z,W,t\bigr),
 \qquad h(t)=1/t.
 \tag{B.4.29.1}
\]
This is a holomorphic automorphism of the
punctured family. The zero section
$e(t)=([1],[1])$ extends to the smooth torus
locus of the central normalization. For
positive integers $j$ tending to infinity,
\[
 t_j=1/j,\qquad t'_j=1/(j+1/2)
\]
lie in the disc and tend to zero, whereas
\[
 T_h(e(t_j))=([1],[1]),\qquad
 T_h(e(t'_j))=([-1],[1]).
\]
In the Tate model these approach the distinct
central smooth points with multiplicative
coordinates $(1,1)$ and $(-1,1)$.
The small resolution is an isomorphism near
both points. Hence $T_h\circ e$ has no
continuous extension to the centre.

This also excludes a meromorphic total
extension, rather than only a holomorphic
one. If such an extension existed, let
$\Gamma\to\mathcal Y$ be its proper
meromorphic graph and pull it back along
$e:\Delta\to\mathcal Y$. The resulting
analytic space is proper over the disc and
has a single point in each noncentral fibre.
Its unique component dominating the disc
is a proper analytic curve, generically of
degree one. After normalization it is finite
of degree one over the normal disc, hence
is the disc itself. Its map to the target
$\mathcal Y$ would extend $T_h\circ e$
holomorphically. This contradicts the two
distinct limits just found. Vertical
components in the pulled-back graph do
not affect this dominating curve.

\emph{Relative forms and their coherent extensions.}
On the ordinary locus the invariant relative
one-forms are
\[
 \omega_1=dZ/Z,\qquad \omega_2=dW/W
 \quad\pmod{dt}.
\]
They and $\omega_1\wedge\omega_2$ are
unchanged by every $T_h$. The holomorphic
forms of a compact torus are invariant, so
this accounts for all relative holomorphic
one- and two-forms on the ordinary fibres.
The relative log forms extend in the Tate
model. Ordinary coherent relative forms
also give finite-order extensions: for
example $t\omega_1,t\omega_2$, and
$t^2\omega_1\wedge\omega_2$ extend as
ordinary relative forms. In the first
resolution chart the first two are explicitly
\[
 t\omega_1=Av\,dx,\qquad
 t\omega_2=Av\,dx+xv\,dA;
\]
the other chart gives the analogous regular
expressions. At an isolated Tate node the
same identity is $t\,dx/x=y\,dx$.

More generally, take the torsion-free image
of each coherent direct image of relative
forms in its bundle on the punctured disc.
Under the two punctured identifications
$\mathrm{id}$ and $T_h$, these are identical
coherent lattices, since the induced map on
the entire ordinary relative-form bundle
is the identity. Coherence therefore gives
no bound on $h$ in this example.

\emph{Gauss--Manin and an absolute form.}
Every fibre translation is homotopic to the
identity and acts identically on integral
cohomology. Consequently $T_h$ induces the
identity on the flat Gauss--Manin bundle,
its connection, and its Hodge subbundle.
The logarithmic period growth and regular
singularity of that connection are unchanged.
They concern the linear cohomology variation,
not the translated origin.

There is even an unchanged absolute form
\[
 \Omega=dt\wedge\frac{dZ}{Z}\wedge\frac{dW}{W}.
 \tag{B.4.29.2}
\]
It descends under both Tate period actions,
since their extra differentials are multiples
of $dt$. In the first small-resolution chart
it is
\[
 d(xAv)\wedge\frac{dx}{x}
          \wedge\frac{d(xA)}{xA}
       =dv\wedge dx\wedge dA,
\]
and in the other chart it is
$dy\wedge dB\wedge du$.
Thus $\Omega$ extends as a nonvanishing
holomorphic canonical form on the smooth
total space. Translation adds
$2\pi i h'(t)dt$ to $dZ/Z$, which vanishes
in (B.4.29.2), so $T_h^*\Omega=\Omega$.

\emph{What a total differential would have to add.}
On a local multiplicative chart, an absolute
lift of the first relative form transforms by
\[
 T_h^*(dZ/Z)=dZ/Z+2\pi i h'(t)dt.
 \tag{B.4.29.3}
\]
This horizontal term is erased both on passing
to relative forms and on wedging with the base
differential. It is not controlled by the
linear Gauss--Manin connection. The stronger
choice $h(t)=e^{1/t}$ makes that term
essential, while leaving all the preceding
relative and canonical calculations unchanged.
For this choice the two sequences
$t_j=1/\log j$ and
$t'_j=1/\log(j+1/2)$ give the same two
distinct central limits. A condition imposing
a regular extending affine connection or
specified absolute lifts could detect such
growth, but it is additional data; properness
and the invariant forms above do not supply it.

The example compares two markings of the same
proper family. It does not identify either
with the actual six-sphere cap. Nor does it
show that differences of two sections both
extending into a single cap fail to extend:
a common essential translation cancels in
such a difference. That stronger possible
bridge remains separate. Exact chart,
monodromy, and translation identities are
checked in
\path{code/e_push_forms_translation.py}.
\end{proof}
